\documentclass[11pt]{letter}
\usepackage[margin=1in]{geometry}
\usepackage{hyperref}

\signature{Elliot Tower \\ Independent researcher \\ elliot@elliottower.ai}
\address{}
\date{\today}

\begin{document}

\begin{letter}{Editorial Office \\ Frontiers in Pharmacology \\ Translational Pharmacology Section}

\opening{Dear Editors,}

We submit ``Drug Failures Partition into Two Kinds: Cross-Design Evidence Diagnosis Across Ten Disease Domains'' for consideration as an Original Research article in Frontiers in Pharmacology.

Nelson et al.\ (2015, Nature Genetics) and Minikel et al.\ (2024, Nature) established that drug targets with genetic support are 2--2.6 times more likely to reach approval. A practical question follows: when MR-based screening misclassifies---predicting failure for an approved drug, or success for a failed one---can the type of error be diagnosed before a pipeline decision is made?

We assembled 41 mechanism families across ten disease domains and applied a deterministic cross-design classification rule comparing observational, Mendelian randomization, and RCT evidence. Three findings address the translational gap between genetic target validation and clinical trial outcomes:

\begin{enumerate}
\item \textbf{Invariance result.} MR-only and the cross-design rule produce identical classifications with zero McNemar disagreements at every threshold tested ($d = 0.05$ to $0.20$). The observational leg carries no predictive information---a structural consequence of the fact that mechanisms reaching Phase~III have non-trivial observational support by construction.

\item \textbf{Two-way failure partition.} The eight misclassified families split by MR effect size into two structurally distinct groups: five with null MR (drug succeeded despite null signal---effector-neutralization and mechanism-bypass) and three with causal MR (drug failed despite a real mechanism---translation gaps and exposure mismatches). Each group implies a different pipeline response.

\item \textbf{Pre-registered design with honest provenance.} Twenty-two families were pre-registered; fourteen were added in two blind amendments (SHAs provided). A provenance table records which boundary classes were pre-specified vs.\ post-hoc. Extension accuracy (6/10, 60\%) is reported separately from pre-registered accuracy (18/22, 82\%) and does not reach significance in isolation.
\end{enumerate}

The paper extends the genetic-support-predicts-approval literature by diagnosing \emph{why} unsupported targets fail and \emph{why} supported targets sometimes fail too---questions that binary genetic-support classification cannot address. The five named boundary classes (effector-neutralization, small-effect, translation-gap, exposure-mismatch, mechanism-bypass) identify where MR screening is expected to produce systematic errors, enabling developers to flag these cases before committing to Phase~III.

\textbf{Suggested reviewers:}
\begin{itemize}
\item Dipender Gill (Imperial College London) --- drug-target MR methodology
\item Iyas Daghlas (University of California, San Francisco) --- MR applications in clinical development
\item Michael V.\ Holmes (University of Oxford) --- MR for drug target validation, pharmacoepidemiology
\end{itemize}

\textbf{Data availability:} The deterministic classifier and all robustness analysis code are provided as supplementary material. All effect sizes derive from published meta-analyses and MR studies cited in the manuscript.

\textbf{Conflict of interest:} The author declares no conflicts of interest.

\closing{Sincerely,}

\end{letter}
\end{document}
